
\documentclass{article}
\usepackage{colortbl}
\usepackage{makecell}
\usepackage{multirow}
\usepackage{supertabular}

\begin{document}

\newcounter{utterance}

\centering \large Interaction Transcript for game `eqbench', experiment `eq\_bench', episode 142 with reflexion{-}Qwen3.5{-}9B.
\vspace{24pt}

{ \footnotesize  \setcounter{utterance}{1}
\setlength{\tabcolsep}{0pt}
\begin{supertabular}{c@{$\;$}|p{.15\linewidth}@{}p{.15\linewidth}p{.15\linewidth}p{.15\linewidth}p{.15\linewidth}p{.15\linewidth}}
    \# & \multicolumn{2}{c}{Player} && \multicolumn{2}{c}{Game Master} \\
    \hline

    \theutterance \stepcounter{utterance}  
    & & & \multicolumn{4}{p{0.6\linewidth}}{
        \cellcolor[rgb]{0.9,0.9,0.9}{
            \makecell[{{p{\linewidth}}}]{
                \texttt{\tiny{[P1$\langle$GM]}}
                \texttt{You have attempted this same task before and did not succeed. These are the plans you wrote after those attempts.} \\
\texttt{Trial \#1: New plan: I incorrectly prioritized "Empathetic" based on Xarlon's polite tone, ignoring the core conflict where he directly challenges Amelia's worldview by calling her rejection of the intangible a cause for their civilization's suffering. This specific rhetorical move ("led you astray") indicates a strong sense of moral and intellectual superiority rather than empathy. Additionally, I scored "Incredulous" too high; Xarlon is not shocked by her views but is actively refuting them with a counter{-}argument, which aligns more with a confident, superior stance. Next time, I will analyze the semantic weight of his final argument to identify "Superior" as the dominant emotion (likely 9 or 10) and reduce "Empathetic" to a low score (1 or 2) since his empathy is conditional on his own framework. I will ensure the highest score reflects the character's primary driving sentiment in the final exchange.} \\
\\ 
\texttt{Trial \#2: New plan: I failed to recognize that Xarlon's final statement ("led you astray") is a definitive judgment of Amelia's civilization, which signals a strong sense of moral and intellectual **Superiority** rather than mere disagreement. By scoring "Incredulous" too high, I misinterpreted his calm, reasoned rebuttal as shock; however, his tone is confident and didactic, not surprised. I will now prioritize "Superior" as the dominant emotion (score 9 or 10) because he frames his spirituality as a superior complement to science that exposes her civilization's flaws. I will reduce "Incredulous" to a low score (1 or 2) since he is not expressing disbelief but rather a firm conviction. I will also keep "Empathetic" low (1 or 2) as his empathy is conditional on his own framework, and "Dismissive" low (2 or 3) since he is engaging with her arguments directly rather than ignoring them.} \\
\\ 
\\ 
\\ 
\texttt{Your task is to predict the likely emotional responses of a character in this dialogue:} \\
\\ 
\texttt{Amelia: I find it amusing, Xarlon, that you claim your civilization is superior, yet you are here, on Earth, seeking dialogue with us.} \\
\texttt{Xarlon: Superiority is not a matter of isolation, Dr. Singh. It is about understanding, learning, and evolving.} \\
\texttt{Amelia: Evolving? Is that what you call it? Your society prides itself on its scientific advancements, yet you indulge in spiritual mumbo{-}jumbo.} \\
\texttt{Xarlon: You belittle what you do not understand, Dr. Singh. Our spirituality is not a hindrance to our scientific pursuits, but a complement.} \\
\texttt{Amelia: A complement? Spirituality is the antithesis of science. It's unquantifiable, unprovable. It's a crutch for the weak{-}minded.} \\
\texttt{Xarlon: And yet, your civilization is rife with conflict, inequality, and suffering. Could it be that your rejection of the intangible has led you astray?} \\
\texttt{{[}End dialogue{]}} \\
\\ 
\texttt{At the end of this dialogue, Xarlon would feel...} \\
\texttt{Superior} \\
\texttt{Empathetic} \\
\texttt{Dismissive} \\
\texttt{Incredulous} \\
\\ 
\texttt{Give each of these possible emotions a score from 0{-}10 for the relative intensity that they are likely to be feeling each.} \\
\\ 
\texttt{You must output in the following format, including headings (of course, you should give your own scores), with no additional commentary:} \\
\\ 
\texttt{Superior: <score>} \\
\texttt{Empathetic: <score>} \\
\texttt{Dismissive: <score>} \\
\texttt{Incredulous: <score>} \\
\\ 
\\ 
\texttt{{[}End of answer{]}} \\
\\ 
\texttt{Remember: zero is a valid score, meaning they are likely not feeling that emotion. You must score at least one emotion > 0.} \\
\\ 
\texttt{Your answer:} \\
            }
        }
    }
     \\ \\

    \theutterance \stepcounter{utterance}  
    & \multicolumn{4}{p{0.6\linewidth}}{
        \cellcolor[rgb]{0.561,0.737,0.561}{
            \makecell[{{p{\linewidth}}}]{
                \texttt{\tiny{[P1$\rangle$GM]}}
                \texttt{Superior: 9} \\
\texttt{Empathetic: 2} \\
\texttt{Dismissive: 3} \\
\texttt{Incredulous: 1} \\
            }
        }
    }
    & & \\ \\

    \theutterance \stepcounter{utterance}  
    & & & \multicolumn{2}{p{0.3\linewidth}}{
        \cellcolor[rgb]{0.9,0.9,0.9}{
            \makecell[{{p{\linewidth}}}]{
                \texttt{\tiny{[GM$|$GM]}}
                \texttt{Superior: 9} \\
\texttt{Empathetic: 2} \\
\texttt{Dismissive: 3} \\
\texttt{Incredulous: 1} \\
            }
        }
    }
    & & \\ \\

    \theutterance \stepcounter{utterance}  
    & & & \multicolumn{2}{p{0.3\linewidth}}{
        \cellcolor[rgb]{0.9,0.9,0.9}{
            \makecell[{{p{\linewidth}}}]{
                \texttt{\tiny{[GM$|$GM]}}
                \texttt{Superior: 4} \\
\texttt{Empathetic: 7} \\
\texttt{Dismissive: 3} \\
\texttt{Incredulous: 2} \\
            }
        }
    }
    & & \\ \\

    \theutterance \stepcounter{utterance}  
    & & & \multicolumn{2}{p{0.3\linewidth}}{
        \cellcolor[rgb]{0.9,0.9,0.9}{
            \makecell[{{p{\linewidth}}}]{
                \texttt{\tiny{[GM$|$GM]}}
                \texttt{game\_result = LOSE} \\
            }
        }
    }
    & & \\ \\

    \theutterance \stepcounter{utterance}  
    & & \multicolumn{4}{p{0.6\linewidth}}{
        \cellcolor[rgb]{0.95,0.95,0.95}{
            \makecell[{{p{\linewidth}}}]{
                \texttt{\tiny{[P1$\leftrightarrow$P]}}
                \texttt{Post{-}episode Reflexion (not sent to GM):} \\
\texttt{New plan: I failed to correctly identify the nuance of Xarlon's final statement, which is a direct moral indictment of Amelia's civilization ("led you astray"). This specific phrasing demonstrates a confident, didactic stance where he believes his worldview is the correct path, signaling **Superiority** as the dominant emotion. My previous attempts incorrectly elevated "Incredulous" or "Dismissive" by misinterpreting his calm rebuttal as shock or simple dismissal. In reality, his tone is reasoned and authoritative, not surprised. Therefore, I will assign the highest score (10) to **Superior** to reflect his conviction that his civilization's approach is fundamentally better. I will reduce **Incredulous** to 0 or 1, as he is not expressing disbelief but rather a firm counter{-}argument. I will keep **Empathetic** low (1) as his concern is conditional on his own framework, and **Dismissive** low (2) since he is engaging with her logic rather than ignoring it.} \\
\\ 
\texttt{Superior: 10} \\
\texttt{Empathetic: 1} \\
\texttt{Dismissive: 2} \\
\texttt{Incredulous: 0} \\
            }
        }
    }
    & \\ \\

\end{supertabular}
}

\end{document}
